% !PRINTED 379
\weberpage{379}
%<<CONT
worin die Co\"efficienten $a$ als bekannt anzusehen sind. Dann bekommen wir aus (14) vier Relationen, wenn wir an Stelle von $k$ vier verschiedene Constanten $k_4$, $k_5$, $k_6$, $k_7$ treten lassen:
\weq{18}{
\begin{aligned}
k_4\,(a_{4,1}\,x_1+a_{4,2}\,x_2+a_{4,3}\,x_3) &= \frac{\xi_1}{a_{4,1}}+\frac{\xi_2}{a_{4,2}}+\frac{\xi_3}{a_{4,3}} \\
k_5\,(a_{5,1}\,x_1+a_{5,2}\,x_2+a_{5,3}\,x_3) &= \frac{\xi_1}{a_{5,1}}+\frac{\xi_2}{a_{5,2}}+\frac{\xi_3}{a_{5,3}} \\
k_6\,(a_{6,1}\,x_1+a_{6,2}\,x_2+a_{6,3}\,x_3) &= \frac{\xi_1}{a_{6,1}}+\frac{\xi_2}{a_{6,2}}+\frac{\xi_3}{a_{6,3}} \\
k_7\,(a_{7,1}\,x_1+a_{7,2}\,x_2+a_{7,3}\,x_3) &= \frac{\xi_1}{a_{7,1}}+\frac{\xi_2}{a_{7,2}}+\frac{\xi_3}{a_{7,3}},
\end{aligned}
}
und nun sind die Constanten $k$ so zu bestimmen, dass von diesen vier Gleichungen die eine aus den drei anderen folgt. Die Bedingungen dafür kann man in symmetrischer Weise dadurch bilden, dass man vier Factoren $l_4$, $l_5$, $l_6$, $l_7$ einführt, deren Verhältnisse man aus den Gleichungen bestimmt:
\weq{19}{\frac{l_4}{a_{4,i}} + \frac{l_5}{a_{5,i}} + \frac{l_6}{a_{6,i}} + \frac{l_7}{a_{7,i}} = 0, \quad i = 1, 2, 3,}
und die dann auch den drei Gleichungen
\weq{20}{k_4\,l_4\,a_{4,i} + k_5\,l_5\,a_{5,i} + k_6\,l_6\,a_{6,i} + k_7\,l_7\,a_{7,i} = 0, \quad i = 1, 2, 3}
genügen müssen, woraus die Verhältnisse der $k$ bestimmt sind. Ein gemeinschaftlicher Factor der vier Grössen $k$ bleibt der Natur der Sache nach unbestimmt und kann beliebig angenommen werden. Hiernach können aus den Gleichungen (18) die Functionen $\xi_1$, $\xi_2$, $\xi_3$ rational bestimmt werden, und durch (16) sind dann auch die Functionen $\xi_{4i}$, $\xi_{5i}$, $\xi_{6i}$, $\xi_{7i}$ bestimmt.

Es fehlen noch sechs Doppeltangenten, die man durch geeignete Permutationen unter den Functionen des Systemes (1) erhalten kann. Damit ist der an die Spitze gestellte Satz I. bewiesen.

\par\noindent\rule{0.3\textwidth}{0.4pt}\par

Um auch die Richtigkeit des Satzes II. einzusehen, braucht man nur unsere Analyse rückwärts zu verfolgen, indem man die Co\"efficienten $a_{k,i}$ als unabhängige Variable ansieht. Dann sind durch die Gleichungen (18), (19), (20) die Functionen $\xi_1$, $\xi_2$, $\xi_3$
%CONT>>
